% Methodology and Artefact (e.g., experiment, framework, prototype etc.) (25 marks)
% - Was there clear description of methodology necessary to reach objectives and fitness for
% purpose (in terms of stated objectives) of the adopted
% methodology/approach/technologies/techniques/etc.;
% - Explain key decisions and justify the choice of methodology and/or tools used;
% - What was the extent to which the objectives were met;
% - Academic/technical challenge of artefact;
% - Was there demonstrable quality in the artefact (depending on the nature of the artefact);
% - Does the artefact provide enough detail to permit reproducibility (incl. datasets, software
% used, a clear methodology, etc.).

% Schema
% - How I did what I wanted to do is the focus
% - What I want to do is: benchmark between heuristics, AI (with my twists) and Hybrid approaches on Dynamic Solomon instances
% To say we have, at the very least:
% - what dynamic means, (customer disclosure over time),
% what stays the same, time windows (hard or soft -> soft in our case), when re-optimisation happens
% what is lateness, what is the optimisation objective (lexicographic, fewer vehicles the better)
% - which heuristics, why, and how I configured them (PyVRP, ORTools, for now)
% - AI -> which AI, how I trained it, on which hardware, with which parameters. With what architecture,
% - AI methodological variants -> which tweaks I applied, why, how (edits to the envs, generators, policies)
% - Hybrid approaches -> How did I made the bridge between AI and classic methods works (warm start, how I implemented it in ORTools and possibly PyVRP)
% - Dynamic setting -> how I made scenario dynamics, why I did it that way, how does the simulator work, performances, observability (visualisation)
% solver agnostic, code publicy available, dod, cutoff, seeds. How did I implement it in ORTOols and why it was much easier for Routefinder.
% how snapshots are built during re-optimisation
% - Instance to problem pipeline -> how are parsed, how are normalised and scaled for RouteFinder
% - How did I setup the benchmark? Which decisions were made? Is it seed-dependent? How do I ensure fairness? How much time do I give heuristics? Role of oracle
% - What am I removing? Stochastic, stochastic travel times, heterogeneous fleet, pickup & delivery, real maps, large-scale 1000-node instances
% - Metrics included and why:  total distance,number of vehicles used, lateness count / lateness sum if soft TW are
% used, rejection count if re-optimisation fails or optional customers are allowed, runtime / wall-clock solve time
% gap to oracle or static baseline where relevant, service level / proportion served on time
% % 

% \section{Intro}
%  1. Problem Definition
%     - what dynamic means, (customer disclosure over time),
%     what stays the same, time windows (hard or soft -> soft in our case), when re-optimisation happens
%     what is lateness, what is the optimisation objective (lexicographic, fewer vehicles the better)
%   - static VRPTW vs dynamic DVRPTW
%   - what “dynamic” means in your work: customer
%     disclosure over time, not stochastic travel
%     times
%   - what remains deterministic
%   - whether time windows are hard or soft in each
%     experiment
%   - what the optimisation objective is
%   - what counts as infeasible, rejected, or late

%   2. Research Questions / Experimental Claims

% \section{Benchmark protocol and fairness}
%   - same dynamic scenarios shared across methods
%   - same random seeds across competitors
%   - same time budget per re-optimisation call
%   - same hardware policy for reporting results
%   - whether AI inference time is counted
%   - whether loading/model initialization is
%     excluded
%   - whether hybrids get extra time or share the
%     same budget
%   - whether oracle comparisons are static full-
%     information only, and not directly comparable
%     to online methods

% \section{Solution Pipeline}

%   - how Solomon instances are parsed
%   - how many base instances you use
%   - how many dynamic replicas per base instance
%   - how dod, cutoff, and seed generate a scenario
%   - whether infeasible dynamic scenarios are
%     discarded or regenerated
%   - how scaling / normalization is handled for
%     RouteFinder tensors
%   - how snapshot instances are built during re-
%     optimisation

% \section{Hybrid Mechanism}

%   - what the AI outputs
%   - how routes/actions are converted into a warm
%     start
%   - how ORTools consumes that warm start
%   - whether PyVRP supports equivalent seeding in
%     your implementation
%   - what happens when the AI proposal is
%     infeasible
%   - whether the hybrid is sequential, repair-
%     based, or improvement-based

% \section{Dynamic Simulator as Artefact}
%   9. Dynamic Simulator as Artefact
%   You already have this, but you should frame it
%   as a reusable benchmark artefact:

%   - solver-agnostic design
%   - event-driven execution
%   - snapshot re-optimisation loop
%   - public code / reproducibility
%   - visual inspection tools
%   - why this architecture was chosen over
%     directly embedding dynamics inside each
%     solver
% \section{Out of scope / Limitations}
%  - assumptions
%   - excluded variants
%   - reasons for exclusion

%   For each excluded feature, say why:

%   - stochastic travel times: excluded to isolate
%     disclosure dynamics
%   - heterogeneous fleets: excluded to avoid
%     conflating fleet assignment with routing
%     quality
%   - pickup and delivery: excluded because
%     RouteFinder and heuristics were compared on a
%     common VRPTW-compatible setting
%   - real maps: excluded for controlled
%     reproducibility
%   - 1000-node instances: excluded due to training
%     and evaluation cost, and because Solomon-
%     scale benchmarking is the target

% \paragraph{AI Setup}
% %Description of the model adopted, RouteFinder.
% Things to discuss
% - training datasets
% - train/validation split logic
% - number of epochs
% - batch size
% - optimizer and learning rate
% - checkpointing / early stopping / model
% selection
% - hardware used
% - inference decoding setup
% - augmentation count at inference
% - whether each model is trained once or with
% multiple seeds
% When training our models, we varied:
% \begin{itemize}
%     \item the number of customers the model has been trained for (50 or 75);
%     \item whether lateness is allowed;
%     \item the geometry of the training set, namely Solomon-based instances, fixed-city
%           locations, or instances generated randomly in the unit square;
%     \item whether the model is trained on dynamic environments or static ones.
% \end{itemize}

% This gives us a total of 8 model families. Every line considers both the 50-customer and
% the 75-customer case.
% \begin{itemize}
%     \item no lateness, Solomon-based, static variants;
%     \item no lateness, general, static variants;
%     \item lateness allowed, Solomon-based, static variants;
%     \item lateness allowed, general, static variants;
%     \item no lateness, general, dynamic variants;
%     \item lateness allowed, general, dynamic variants;
%     \item no lateness, Solomon-based, dynamic variants;
%     \item lateness allowed, Solomon-based, dynamic variants.
% \end{itemize}

% % description of how I varied the various envs and why
% To do so, we customised the base env found in RouteFinder to edit,
% (for each variant)
% - number of customers -> not edit to the code, just a different parameter passed
% - lateness -> changed the env (env_flexible.py)
% - geometry -> Solomon_generator -> the generator now takes in input a list of Solomon instances,
% and based on a counter gives one or another Solomon instance when generating a batch. The
% sampling of locations has been disabled.
% - dynamic training -> see below
% \paragraph{OR Tools setup}
% % description of how I set up ortools and why I did it the way I did.

% \paragraph{Hybrid Setup}
% % how I setup ortools to work with a warm start

% The experiments were run on three types of hardware: low-end, mid-range, and high-end.
%dynamic simulator section

% \paragraph{Dynamic Setup}
% The dynamic benchmark was implemented as an event-driven simulator that converts each
% static Solomon instance into a partially revealed DVRPTW scenario before execution.
% For a given degree of dynamism, $\epsilon$, the simulator samples $\mathrm{round}(\epsilon N)$
% customers to be dynamic, where $N$ is the total number of customers. Each selected
% customer is assigned a reveal time drawn uniformly from the interval
% $[0, \min(l_i, \tau)]$, where $l_i$ is the customer's original due time and
% $\tau = \text{cutoff\_ratio} \times T$, with $T$ denoting the depot horizon. Once
% revealed, the customer becomes available only from that reveal instant onward,
% implemented by updating its ready time to $\max(e_i, r_i)$, where $e_i$ is the
% original ready time and $r_i$ is the reveal time; due times remain bounded by the
% original planning horizon.

% During simulation, vehicles are advanced exactly to each reveal event while preserving
% their current position, elapsed time, residual capacity, and any non-preemptive service
% already in progress. At every event, the simulator constructs a residual snapshot
% instance containing all currently active but unserved customers together with one
% artificial start node per vehicle at its current state, and the chosen solver is invoked
% to re-optimise the remaining routes under the allotted time budget. If re-optimisation
% fails, the newly revealed request is rejected and the previous plan is retained. Final
% performance is measured on the executed routes produced by this sequential
% execution-reoptimisation process, with feasibility assessed using unserved customers,
% time-window violations, and capacity violations.

% As long as it takes as input an object of type \texttt{Instance} and returns a
% \texttt{Solution}, every solver can be used. This has been fundamental in enabling us to
% use different solvers and models to conduct our comparative study.

% The solver-agnostic simulator allowed us to benchmark different approaches seamlessly.
% For the metaheuristics part, we used ORTools and PyVRP. For the AI component, we took
% the RouteFinder \cite{RouteFinder} code and trained it in different ways to understand
% the impact of those changes on the final result.
%end of dynamic simulator section

% lateness section
% A key modification required to support lateness
%   in the RouteFinder-based model was the
%   introduction of a custom environment,
%   MTVRPFlexibleEnv, implemented in src/
%   dvrptw_bench/rl/env_flexible.py:100. The
%   standard RouteFinder MTVRPEnv treats customer
%   time windows as hard constraints: if a vehicle
%   cannot begin service before a customer’s due
%   time, that customer is masked out and cannot be
%   selected. This behaviour is unsuitable for the
%   soft time-window setting used in this work,
%   where late service is permitted but should
%   incur a penalty. The flexible environment
%   therefore extends the original formulation by
%   introducing configuration parameters such as
%   allow_late_customers, lateness_penalty, and
%   hard_depot_time_window. These parameters allow
%   customer deadlines to be relaxed while
%   optionally preserving the depot closing time as
%   a hard constraint.

%   The environment dynamics were modified in four
%   main ways. First, during each transition, the
%   environment now computes the service start time
%   at the selected customer and accumulates
%   tardiness as the positive difference between
%   service start and due time. This value is
%   stored as total_tardiness in the state
%   representation and updated incrementally
%   throughout the route construction process.
%   Second, the action mask was altered so that,
%   when lateness is enabled, customers are no
%   longer excluded solely because they would be
%   reached after their due time. In other words,
%   customer time windows become soft at decoding
%   time, whereas the depot horizon can still
%   remain hard if required. Third, the reward
%   function was extended from pure tour length
%   minimisation to a penalised objective that
%   combines travelled distance with a lateness
%   cost proportional to accumulated tardiness.
%   This ensures that the model does not simply
%   exploit soft feasibility, but instead learns to
%   trade off route efficiency against the cost of
%   serving customers late. Fourth, the solution
%   validity checks were relaxed consistently with
%   this new formulation, so that customer lateness
%   is no longer treated as infeasibility when
%   allow_late_customers=True, while depot-deadline
%   feasibility can still be enforced separately.

%   No structural changes were required to the
%   RouteFinder policy itself. The original
%   RouteFinderPolicy was retained unchanged, since
%   it already embeds the information needed to
%   reason about time-window decisions, including
%   customer ready times, due times, service
%   durations, and the current route time.
%   Likewise, no changes were required to the
%   standard MTVRPGenerator. The generator already
%   produces conventional VRPTW instances with
%   coordinates, demands, service times, and time
%   windows; the distinction between hard and soft
%   time windows is therefore not a property of the
%   generated data, but of how the environment
%   interprets those data during training and
%   inference. Consequently, support for lateness
%   was achieved entirely through environment-level
%   modifications rather than through changes to
%   the neural architecture or the instance
%   generator.
% end lateness section

% start of dynamic training section
% • To train RouteFinder on dynamic instances
% rather than static VRPTW instances, it was
% necessary to modify both the instance
% generation process and the environment
% dynamics. This was implemented through the
% DynamicGenerator in src/dvrptw_bench/
% paper_dynamic_RouteFinder/generator.py:1 and
% the MTVRPDynamicEnv in src/dvrptw_bench/rl/
% env_dynamic.py:1. The objective of these
% modifications was to expose the model, during
% training, to the same partial-information
% setting that characterises the dynamic DVRPTW:
% at the start of an episode only a subset of
% customers is known, while the remaining
% requests are revealed progressively over time.

% For each sampled instance, it selects a proportion
% of customers to become dynamic according to a
% degree of dynamism parameter. Each selected
% customer is assigned a reveal time drawn
% uniformly between zero and an upper bound
% determined by both the customer’s due time and
% a disclosure cutoff ratio applied to the depot
% horizon. Once a reveal time has been assigned,
% the customer’s ready time is updated to the
% maximum of its original ready time and its
% reveal time. In this way, a customer cannot be
% served before it is disclosed, while still
% preserving the original time-window structure
% as much as possible. The generator then returns
% not only the usual routing features such as
% coordinates, demands, service times, and time
% windows, but also dynamic metadata, in
% particular reveal_times, an is_dynamic
% indicator, and the corresponding disclosure
% schedule. This effectively transforms each
% static Solomon instance into a partially
% revealed training scenario.

% The dynamic environment is responsible for
% making those disclosure times operational
% during decoding. Instead of exposing the full
% customer set from the beginning, the
% environment maintains both the complete
% underlying instance and the currently visible
% version presented to the policy. This is
% achieved by storing two parallel
% representations: the “actual” tensors, which
% contain the full problem information, and the
% observation tensors, which hide customers that
% have not yet been revealed. A Boolean revealed
% mask is updated throughout the episode by
% comparing each customer’s reveal time with the
% elapsed simulation time. When hidden-customer
% masking is enabled, unrevealed customers are
% replaced in the observation by depot-level
% placeholders with zero demand and neutral
% service information, ensuring that the policy
% can only act on currently disclosed requests.

% A further change required for dynamic training
% was the introduction of an explicit notion of
% elapsed time distinct from route time. The
% environment tracks the physical progression of
% time through an elapsed_time variable, which
% increases as the vehicle travels, waits, and
% serves customers. After each decision, the
% environment checks whether any unrevealed
% customers have now become available. If no
% currently revealed customer can be feasibly
% visited, but some future customer remains
% undisclosed, the environment automatically
% advances time to the next reveal event. This
% waiting mechanism is essential: it allows the
% episode to continue realistically even when the
% vehicle has no immediate feasible action among
% the customers currently known. As a result,
% training episodes reproduce the event-driven
% nature of the dynamic problem, where routing
% decisions are repeatedly revised as new
% requests appear.

% Feasibility checking was also adapted to the
% dynamic setting. In the static RouteFinder
% environment, feasibility is computed over the
% full customer set. In the dynamic version,
% feasibility is still evaluated using the actual
% customer attributes, but a customer is only
% selectable if it is both feasible and already
% revealed. Thus, the action mask enforces
% partial observability directly: unrevealed
% customers cannot be chosen even if they would
% otherwise satisfy capacity and time-window
% constraints. The reward function itself remains
% based on travelled distance, but the sequence
% of decisions is now conditioned on progressive
% disclosure rather than full information from
% the outset.

% Together, these components allow
% RouteFinder to be trained on trajectories that
% more closely match the true online re-
% optimisation setting considered in this work,
% rather than on conventional static VRPTW
% episodes.

%end of dynamic training section

\chapter{Methodology}
\label{chp:methodology}

This chapter describes the experimental framework used to compare classical
optimisation, AI-based, and hybrid solution methods for the Dynamic Vehicle
Routing Problem with Time Windows (DVRPTW). The study focuses on dynamic
deterministic customer disclosure scenarios derived from the Solomon benchmark
set. All 56 benchmark instances are considered, spanning the three standard
Solomon families (C, R, and RC) and including both Type 1 and Type 2 variants.
For each instance, a subset of the 100 customers will be used for reasons better
explained in the limitations section\ref{limitations}.


The chapter first defines the problem setting and benchmark construction
process, followed by the configuration of the compared methods. It then presents
the dynamic simulation environment, fairness protocol, and evaluation metrics
used to assess solution quality, computational efficiency, and service
performance.

To support reproducibility, the benchmark artefact is intended to be released as
open-source code in a public repository [GitHub repository placeholder:
        \texttt{https://github.com/username/dvrptw-benchmark}]. The repository is
intended to include the simulator, data-processing pipeline, solver adapters,
training scripts, benchmark configuration files, and visualisation utilities
used in this study. Together with the Solomon benchmark instances, fixed random
seeds, and the reported experimental settings, this should allow the benchmark
scenarios, training procedures, and evaluation runs to be reproduced by other
researchers.


\section{Experimental Problem Setting}
{\color{red} discuss the hardware used for evaluation (low,mid,high-end) in
evaluation}

In this study, dynamicity is modelled through the progressive disclosure of a
subset of customers from the Solomon benchmark instances. Progressive customer
disclosure is the only source of dynamicity considered. Customer locations,
demands, travel times, service times, and time windows remain deterministic and
fixed once a scenario has been generated. Consequently, the resulting benchmark
belongs to the class of dynamic deterministic vehicle routing problems.
\cite{ojedariosRecentDynamicVehicle2021}

Both hard and soft time-window settings are considered. Under hard time windows,
service must begin within the customer interval. Under soft time windows, late
service is permitted but incurs a penalty. This allows comparison between
methods under both strict feasibility requirements and more flexible service
regimes.

Re-optimisation occurs whenever a new customer request is revealed. At each
event, the current execution state is converted into a residual snapshot
instance, which is then re-solved from that point onward using the selected
method.

\subsection*{Degree of Dynamism}

The degree of dynamism (DoD), denoted by $\epsilon$, controls the proportion of
customers that are unavailable at time zero and are instead revealed during
execution. For an instance containing $N$ customers, $\mathrm{round}(\epsilon
N)$ customers are designated as dynamic requests.

The benchmark evaluates the following levels:

\[
    \epsilon \in \{0.10, 0.25, 0.50, 0.75\}.
\]

These values represent progressively more challenging scenarios, ranging from
low to high dynamicity.

\subsection*{Reveal Horizon}

For each dynamic customer, the reveal time is sampled uniformly from

\[
[0,\tau], \qquad \tau = \alpha T,
\]

where $T$ is the depot planning horizon and $\alpha = 0.8$. Hence, all dynamic
requests are disclosed within the first 80\% of the planning horizon.

This avoids trivial late disclosures while preserving substantial online
uncertainty during execution.

\section{Simulator}
The dynamic benchmark was implemented as an event-driven simulator that converts
each static Solomon instance into a partially revealed DVRPTW scenario prior to
execution. For a given degree of dynamism, $\epsilon$, the simulator samples
$\mathrm{round}(\epsilon N)$ customers to become dynamic, where $N$ is the total
number of customers. Each selected customer is assigned a reveal time drawn
uniformly from the interval $[0,
        \min(l_i,\tau)]$,
where $l_i$ is the original due time and $\tau=\text{cutoff\_ratio}\times T$,
with $T$ denoting the depot planning horizon. Once revealed, the customer
becomes available only from that instant onward by updating its ready time to
$\max(e_i,r_i)$, where $e_i$ is the original ready time and $r_i$ is the sampled
reveal time. Original due times are preserved.

During execution, vehicles are advanced exactly to each reveal event while
preserving their current position, elapsed time, residual capacity, travelled
distance, and any non-preemptive service already in progress. At each event, the
simulator constructs a residual snapshot instance containing all currently
revealed unserved customers together with one artificial start node per vehicle
representing its current execution state. The selected solver is then invoked to
re-optimise the remaining routes under the shared time budget.

If re-optimisation fails to incorporate a newly revealed request, that request
is recorded as rejected and the previously committed plan is retained. This is a
known limitation that will be discussed here\ref{limitations}. Final performance
is evaluated on the routes actually executed over time, rather than on
intermediate planned solutions, with feasibility assessed through unserved
customers, capacity violations, and time-window violations.

The simulator was designed to be solver-agnostic: any method capable of taking
an
\texttt{Instance}
object as input and returning a
\texttt{Solution}
object can be integrated into the benchmark. This common interface was central
to the comparative study, as it enabled consistent evaluation of heuristic,
AI-based, and hybrid approaches within the same dynamic execution framework.
\subsection{Simulator Visualisation}
Visualisation was used as a supporting tool for qualitative inspection of both
static and dynamic solutions. In particular, it was used to verify the
correctness of generated scenarios, to inspect the timing of customer disclosure
events, and to observe how routes evolved across re- optimisation steps in the
dynamic setting. The visualisation component also includes an interactive
inspection capability, which was useful for debugging individual instances by
allowing closer examination of customer visibility, vehicle progression, route
changes, and the effect of successive reveal events. This was especially
valuable for validating simulator behaviour and for interpreting differences
between heuristic, AI, and hybrid solution strategies, although visual analysis
was not used as a primary benchmark metric. {\color{red}Simulator image}



\section{Benchmark Pipeline and Solution Methods}

{\color{red} Remember a pipeline image}

This section describes the three families of solution methods compared in the
benchmark: classical heuristics, AI-based policies, and hybrid approaches. Since
all methods are evaluated on the same dynamic Solomon scenarios, the section
first defines the common representation pipeline used to parse raw benchmark
instances, transform them into solver-ready problem objects, and reconstruct
residual snapshots during online re-optimisation. This shared pipeline ensures
that performance differences are attributable to the solution methods rather
than to inconsistent input processing or problem encoding.

\paragraph{Common Instance Representation}
The raw Solomon file is first parsed into a structured VRPTW instance by
extracting the depot, the list of customers, the available number of vehicles,
the vehicle capacity, and the time-window and service information associated
with each node. Each customer is therefore represented internally by its
coordinates, demand, ready time, due time, and service duration. After this, a
full distance matrix is computed over the depot and customer set, producing a
common VRPTWInstance representation. This intermediate format is
solver-independent and serves as the shared input object for all methods in the
benchmark, including heuristics, AI models, and hybrid approaches.

\paragraph{Dynamic Snapshot Construction}
Snapshots are constructed as residual problem states whenever a new customer
reveal event occurs. At each event time, the simulator first advances all
vehicles from their previous states up to the current time, updating their
positions, elapsed times, remaining capacities, travelled distances, and any
service already in progress. Customers are then partitioned into those already
served, those newly or previously revealed and still awaiting service, and those
temporarily unavailable due to ongoing non-preemptive service. Only currently
relevant and serviceable customers are retained in the snapshot.

Each snapshot stores the current simulation time together with the active
customer set and the full state of every vehicle. Vehicle states include current
location, residual capacity, elapsed time, route progress, travelled distance,
and any active service commitment. The snapshot therefore represents the
complete residual online routing state rather than only a reduced customer list.

For re-optimisation, the snapshot is transformed into a residual VRPTW instance.
Vehicles that are no longer located at the depot are represented by artificial
start nodes encoding their current positions and earliest feasible departure
times, while idle vehicles may continue to start from the physical depot. A
common depot is retained as the route end point, and a new distance matrix is
computed over the resulting node set.

This produces a multi-start residual routing instance containing all currently
available unserved customers, one start state per vehicle, remaining capacities,
and updated temporal constraints. As a result, each solver re-optimises from the
true execution state at the time of disclosure rather than restarting from the
original static benchmark instance.

\subsection{Classical Heuristics Methods}
\paragraph{ORTools}
The first classical baseline is OR-Tools\cite{ortools_routing}, used here as a
time-bounded routing solver. ORTools is an open-source optimisation toolkit
developed by Google that provides efficient algorithms for solving combinatorial
optimisation problems, including vehicle routing problems. ORTools was used in
this work as a classical heuristic baseline for the VRPTW and DVRPTW
experiments. It was configured as a time-bounded routing solver that combines a
constructive phase with a subsequent improvement phase. The exact configuration
can be found in the evaluation section \ref{chp:evaluation}. First, it generates
an initial feasible solution using a greedy route-building strategy. This
initial plan is then refined through local search, with the objective of
reducing total travel distance while maintaining feasibility with respect to
vehicle capacities and time-related constraints. In this way, ORTools is used
not as an exact optimisation method, but as a practical metaheuristic baseline
designed to produce strong solutions within short computational budgets.

The solver was configured to model the standard VRPTW constraints explicitly.
Vehicle capacities were enforced throughout each route, and travel time was
represented jointly with service duration so that each customer had to be
visited within its allowable service interval. In the static experiments, time
windows were treated as hard constraints. In the dynamic experiments, a soft
time-window configuration was also used, allowing late service when necessary
but attaching a penalty to such violations. This design choice was important
because, in the dynamic setting, new requests may be revealed after routes are
already under execution, making strict time-window feasibility harder to
preserve under short re-optimisation times.

ORTools was executed under the shared runtime policy described in
Section~\ref{sec:fairness}. In dynamic experiments, the solver was invoked on
each residual snapshot using the allocated per-event budget.

In the hybrid experiments, ORTools was also used as the downstream improvement
solver. Its hybrid role is described separately in Section~\ref{sec:hybrid}.

The exact parameter settings used for OR-Tools are reported in the evaluation
and reproducibility sections. Only limited tuning was required, as the adopted
configuration follows standard practice in the OR-Tools routing literature,
using well-established first-solution and local-search strategies designed for
general VRPTW instances.

% The stopping criterion for ORTools was based entirely on a fixed runtime
% budget. In the static benchmark, each instance was solved for at most the
% allocated time limit. In the dynamic benchmark, the same principle was applied
% at every re-optimisation event: whenever new customer information became
% available, the solver was invoked on the residual problem and allowed to search
% only within the assigned per-event budget. This ensured that ORTools remained
% directly comparable with the AI and hybrid methods under the same online
% decision-making conditions.

% In the hybrid setting, ORTools also served as a route-improvement component.
% Instead of always starting from scratch, it could be seeded with an initial plan
% produced by the AI model. The role of ORTools in this case was to refine that
% candidate solution within the same fixed runtime budget, thereby combining the
% constructive guidance of the learned policy with the improvement capability of a
% classical local-search-based optimiser.
\paragraph{PyVRP (still undecided...)}

\subsection{AI Methods}

\subsubsection{RouteFinder (Base Setup)}

RouteFinder was adopted as the primary AI-based routing method in this study. It
was selected due to its strong reported benchmark performance and open-source
implementation, which allowed controlled retraining and modification.
\cite{RouteFinder}

The focus of this methodology is therefore not the original RouteFinder
architecture itself, but how the framework was configured, adapted, and
evaluated within the proposed DVRPTW benchmark.

\paragraph{RouteFinder Pipeline from VRPTW Instance to Solution}
{\color{red} I'm considering moving more low-level details to the appendix here,
such as TensorDict mention, depot tokens as separators, internal node indices.
Mention the fact that we ignore training cost as we are considering an
inference-only scenario} Given the shared
\texttt{VRPTWInstance},
RouteFinder requires an additional conversion step into its tensor-based input
format. The instance is converted into a
\texttt{TensorDict}
containing customer coordinates, normalised demands, service times, time
windows, vehicle capacity, and related routing attributes. When coordinate
normalisation is enabled, spatial coordinates are rescaled to the unit square,
and time-related quantities are scaled consistently so that travel distance and
travel time remain aligned with the unit-speed assumption used by the model.

The resulting tensor representation is then passed to the RouteFinder
environment, which initialises the decoding state. The trained model generates a
route by sequentially selecting feasible actions according to its learned
policy. Inference may be performed greedily or with test-time augmentation. In
the latter case, multiple transformed versions of the same instance are
evaluated, and the best decoded route is retained.

The decoded output is an action sequence rather than a benchmark-level solution.
Therefore, a final reconstruction step maps depot tokens to route separators and
translates RouteFinder's internal customer indices back to the original Solomon
customer identifiers. The resulting vehicle routes are stored in the common
\texttt{Solution} format, allowing the benchmark framework to compute distance,
runtime, feasibility, and service-quality metrics consistently across methods.

\subsubsection{Training Variants}

{\color{red} Probably need a citation to confirm why we would need training
variants at all...} {\color{red} Use Ablation table in evaluation}

To examine how training conditions influence the performance of neural routing
policies in dynamic vehicle routing, several RouteFinder variants were trained
under different data-generation assumptions. These can be grouped into three
main methodological directions: a general synthetic training regime, a
Solomon-shaped training regime, and a dynamic training regime. Each direction
was also studied under different problem sizes and under both strict and
lateness-enabled settings.

\paragraph{General Synthetic Variant}

The first training direction used general synthetic VRPTW instances generated
without explicitly reproducing the structural characteristics of Solomon
benchmarks. Customer locations and associated routing attributes were sampled
more generically, producing a broad distribution of instances intended to
encourage generalisation. This variant therefore serves as the most
distribution-agnostic learned baseline.

\paragraph{Solomon-Shaped Variant}

The second training direction aimed to reduce mismatch between the training
distribution and the benchmark distribution by constructing instances that more
closely resemble Solomon problems. Generated training data reflected geometric
and temporal regularities commonly found in Solomon benchmarks, such as
clustered customer layouts, depot-centred routing patterns, and tighter
benchmark-like time windows. This variant can therefore be interpreted as a
distribution-aligned training strategy.

\paragraph{Dynamic Training Variant}

The third training direction focused on the disclosure process itself. Rather
than training only on fully known static VRPTW instances, this variant trained
the model on partially revealed routing problems in which some customers became
available only during execution. The objective was to align training more
closely with the true DVRPTW setting by exposing the policy to incomplete
information and repeated route revision.

\paragraph{Lateness and Problem-Size Variants}

Across the above training directions, additional variants were introduced by
changing two further factors. First, models were trained for different problem
sizes, specifically 50-customer and 75-customer instances. Second, both strict
and lateness-enabled versions were considered. In the strict setting, time
windows were treated as hard constraints, whereas in the lateness-enabled
setting, late service was permitted but penalised.

\paragraph{Combined Model Families}

Taken together, these design dimensions define a family of trained models rather
than a single learned solver. The experiments therefore compare not only
heuristics, AI, and hybrid methods, but also alternative AI training regimes.

\subsubsection{Training Setup}
{\color{red}I must either try training with multiple seeds or justify why
training variance is negligible.}

All neural models were trained using the RouteFinder architecture within the
same general optimisation framework, while varying the training distribution
according to the methodological variant under study. In all cases, training data
were generated online rather than assembled as a fixed stored dataset.

Consequently, the distinction between training and validation data did not rely
on a conventional file-based split, but on separate generated samples: a large
stream of instances was used for training, while a smaller independently
generated set was used for validation under the same generation regime. This
design ensured that validation measured generalisation within the target
distribution rather than memorisation of a finite benchmark subset.

For most experiments, the training configuration followed a common pattern.
Models were trained for 200 epochs, with a batch size of 256 instances per
optimisation step. The training set size per epoch was fixed at 100{,}000
generated instances, while the validation set contained 10{,}000 instances.

Optimisation was carried out using the default RouteFinder training framework
with Adam-style gradient-based optimisation, a learning rate of $3
    \times 10^{-4}$, and
a weight decay of $10^{-6}$. These settings were kept constant across variants
in order to isolate the effect of the training distribution and environment
design.

Checkpointing was used throughout training to preserve intermediate models and
support resumption after interruption. No early stopping strategy was applied,
and model selection was based primarily on the final trained checkpoint.

At inference time, decoding was performed either greedily or using test-time
augmentation. Under augmented decoding, multiple transformed versions of the
same instance were evaluated and the best resulting solution retained. The
typical augmentation count used at inference was eight.

Finally, each model configuration was trained once rather than across multiple
independent training seeds. Randomness was instead controlled more explicitly at
the benchmark level through the seeds used to generate dynamic scenarios during
evaluation.

\subsection{Hybrid Mechanisms}
\label{sec:hybrid}
{\color{red} I'm considering transforming it into a smoother prose. It came out
as a list of things... Probably needs a citation or two}

The hybrid mechanism follows an improvement-based strategy. RouteFinder first
generates an initial route plan for the current static instance or residual
dynamic snapshot. This plan is reconstructed into the shared
\texttt{Solution} format and then translated into the internal route assignment
required by ORTools. ORTools subsequently performs local search from this
incumbent solution, rather than constructing a solution from scratch.

The motivation is to combine complementary strengths: RouteFinder provides fast
constructive guidance from a learned policy, while ORTools provides mature
neighbourhood search and constraint-handling mechanisms. In the dynamic setting,
this is intended to reduce the search effort required at each re-optimisation
event.

\paragraph{Budget policy}
Each optimisation event is assigned a fixed wall-clock budget. For the hybrid
method, RouteFinder inference time is counted within the same total budget
anddeducted before OR-Tools refinement begins. By contrast, benchmark-management
overheads external to the optimisation process, including snapshot generation
and the translation of AI solutions into the OR-Tools internal assignment
format, were not charged against the runtime budget.
\paragraph{AI Failure Handling}
If the AI-generated proposal cannot be translated into a valid assignment for
the current snapshot, the solver falls back to its standard initialisation
procedure. This ensures that the hybrid method remains robust even when the
learned policy performs poorly on out-of-distribution scenarios.


\section{Benchmark Protocol and Fairness}

To ensure a fair comparison across heuristics, AI, and hybrid approaches, all
methods were evaluated on exactly the same dynamically generated scenarios,
using the same random seeds and the same per-event re- optimisation budget. A
common hardware reporting policy was also maintained so that runtime results
could be interpreted consistently. In the online setting, each solver invocation
was subject to the same time limit, and no method was intentionally granted a
larger re-optimisation budget than another. AI inference time, possibly
including augmentation and sampling, are included in the time budget, whilst
one-off costs such as model loading and initialisations are not. When working in
flexible environment with soft constraints, it's important to assign penalties
to behaviours that violate the constraints defined for the problem. We will now
discuss how we handled such penalties in our benchmark.

\subsection{Penalties}

When soft time-window settings are enabled, solutions are evaluated through a
composite objective that balances routing efficiency against service violations.
This is necessary because allowing lateness without an explicit penalty would
make comparisons across methods ambiguous and could incentivise unrealistic
solutions with excessive delay.

Let $K$ denote the number of vehicles used, $D$ the total travelled distance,
and $L=\sum_i \max(0,a_i-l_i)$ the cumulative lateness, where $a_i$ is the
actual service start time at customer $i$ and $l_i$ is the corresponding due
time. The soft time-window objective is evaluated lexicographically as:

\[
\min \; (K,\; D + \lambda L)
\]

That is, the number of vehicles remains the primary criterion, consistent with
classical Solomon-style evaluation, while travelled distance plus penalised
lateness is minimised as the secondary criterion.

The coefficient $\lambda$ controls the trade-off between route cost and delayed
service. A larger value makes lateness progressively less attractive, thereby
approaching hard time-window behaviour, whereas a smaller value allows greater
schedule flexibility.

For the methods considered in this study, the following values were used:

\begin{itemize}
    \item \textbf{OR-Tools:} $\lambda = 50$
    \item \textbf{RouteFinder-based models:} $\lambda = 150$
    \item \textbf{Hybrid RouteFinder + OR-Tools:} the AI proposal is generated
under the RouteFinder penalty ($\lambda = 150$), after which OR-Tools refinement
uses $\lambda = 50$ for final improvement.
\end{itemize}

The OR-Tools value was selected empirically through preliminary experiments as a
practical compromise between discouraging lateness and preserving search
flexibility. The larger RouteFinder penalty was adopted to provide a stronger
training signal against tardy solutions during learning, where softer penalties
were observed to tolerate excessive lateness. Although the coefficients differ
across paradigms, each value was calibrated to induce similar behavioural
preference toward timely service within the corresponding optimisation process.

No additional penalty was applied for rejected customers in the simulator.
Instead, rejected requests are recorded separately through rejection-count and
service-level metrics, since rejection is treated in this study as a dynamic
service outcome rather than as a direct optimisation term.


\section{Evaluation Metrics}
{\color{red} Citations to dvrp surveys and schyns}

In this study, solver quality cannot be assessed by a single scalar because
route cost, service feasibility, and computational responsiveness are all
materially relevant.
\cite{valueOfInformation, ojedariosRecentDynamicVehicle2021,psaraftisDynamicVehicleRoutingProblems2016}.

All metrics are computed from the final executed solution, rather than from
intermediate planned routes, so that results reflect actual online performance
under progressive customer disclosure.


\subsection{Operational Metrics}

Operational metrics evaluate the efficiency of the produced routing plan.

\begin{itemize}
    \item \textbf{Total distance}: the cumulative distance travelled by all vehicles across the planning horizon. This is the primary proxy for transportation cost and fleet efficiency, and remains one of the most widely reported objectives in VRP studies.

    \item \textbf{Number of vehicles used}: the number of vehicles that are activated to serve at least one customer. This metric is relevant because fleet utilisation often carries fixed operational costs and is commonly prioritised in lexicographic VRPTW objectives such as the Solomon Benchmark.
\end{itemize}

\subsection{Service Quality Metrics}

Service metrics evaluate how effectively customer demand is satisfied under
dynamic conditions.

\begin{itemize}
    \item \textbf{Lateness count}: the number of customers served after the end of their time window. This indicates how frequently service commitments are violated.

    \item \textbf{Lateness sum}: the cumulative delay across all late customers, typically measured as the sum of positive tardiness values. This captures the severity of service failures rather than only their frequency.

    \item \textbf{Rejection count}: the number of customer requests that remain unserved or are explicitly discarded. In dynamic settings, this reflects the capacity of a method to absorb newly revealed demand.

    \item \textbf{Service level}: the proportion of revealed customers out of all customers successfully served within their time windows. This provides an interpretable customer-oriented measure of reliability.
    \item \textbf{Responsivness}: {\color{red} Wanted to add it, but unsure about how to measure it. I'll check \cite{AntColonySchyns2015} again.}
\end{itemize}


\subsection{Computational Metrics}

Computational metrics assess whether each method can operate effectively within
the real-time decision limits of the benchmark. Since all solvers are evaluated
under the same runtime cap, these measures do not reflect unrestricted solver
speed, but rather how efficiently each method uses the available computation
budget and whether high-quality decisions can be produced within the allowed
time.

\begin{itemize}
    \item \textbf{Solve time}: the wall-clock time required for a solver call on a static instance or dynamic snapshot.

    \item \textbf{Cumulative runtime}: the total wall-clock time consumed across all re-optimisation events during a dynamic scenario.
\end{itemize}

\subsection{Comparative Metrics}

Comparative metrics position online performance relative to stronger-information
baselines, such as the oracle.
\paragraph{Oracle}

The oracle baseline is included as a full-information reference used to
contextualise the performance of online methods. It is not treated as a direct
competitor, since unlike the dynamic approaches it has complete advance
knowledge of all customer requests, their release times, and all relevant
instance information from the start of the planning horizon.

Operationally, the oracle is obtained by solving the corresponding static
full-information instance using OR-Tools with Guided Local Search (GLS) as the
improvement metaheuristic. A runtime budget of five minutes is allocated per
instance, substantially larger than the online per-event budgets used in the
dynamic experiments, in order to obtain a strong high-quality reference
solution. The oracle objective is the minimisation of route cost, subject to the
relevant problem constraints.

Because the oracle is itself heuristic rather than provably optimal, it should
be interpreted as a strong practical upper benchmark rather than a guaranteed
global optimum. Accordingly, gaps to oracle measure the penalty of operating
under incomplete information relative to a strong full-information solution, not
the distance to mathematical optimality. From the oracle solution we are able to
compute the value of information.

\begin{itemize}
    \item \textbf{Value of information}: the improvement achievable when future request information is known in advance. This metric is frequently used in dynamic routing studies to quantify the operational benefit of anticipatory information \cite{valueOfInformation}.
\end{itemize}

\subsection{Metric Interpretation for This Study}

Because the study compares heuristic, AI-based, and hybrid methods, the metrics
also support different research questions. Operational metrics indicate route
quality, service metrics capture robustness under dynamic arrivals,
computational metrics assess practical deployability, and comparative metrics
help quantify the penalty of online decision-making. Together, these measures
provide a balanced basis for evaluating whether learned policies or hybrid
warm-start approaches offer advantages over classical optimisation baselines.
\section{Reproducibility Statement}
{\color{red} Should I put those into evaluation (or even appendix?). I plan to
add tables with versions, seeds, etc etc.}



\section{Limitations and Out of Scope}
\label{limitations}

In order to maintain a controlled and interpretable experimental setting,
several important extensions of real-world routing were intentionally excluded
from the present study. These exclusions help isolate the central research
question, namely the comparison of heuristic, AI-based, and hybrid methods under
dynamic customer disclosure, but they also define the scope and limitations of
the reported results.

\begin{itemize}

    \item \textbf{Stochastic travel times.}
Travel times were assumed deterministic and known in advance. Random congestion,
incidents, and time-dependent traffic effects were not modelled. This
simplification was adopted to isolate the impact of customer disclosure dynamics
without conflating results with travel-time uncertainty.

    \item \textbf{Heterogeneous fleets.}
All vehicles were assumed identical in capacity and operating characteristics.
Variations in fleet composition, such as mixed capacities or vehicle-specific
costs, were excluded to avoid combining fleet assignment decisions with
routing-quality comparisons.

    \item \textbf{Pickup and delivery constraints.}
The benchmark considers standard delivery-style VRPTW requests only. Paired
pickup-and-delivery operations, precedence constraints, and load rebalancing
requirements were not included, allowing all compared methods to operate within
a common VRPTW-compatible setting.

    \item \textbf{Real road networks and map data.}
Distances were derived from benchmark coordinates rather than from real road
networks. As a result, factors such as asymmetric travel times, turn
restrictions, and urban traffic structure were not considered. This choice was
made to preserve controlled reproducibility and comparability with prior
Solomon-based studies.

    \item \textbf{Large-scale instances.}
Large-scale routing instances involving hundreds or thousands of customers were
outside the scope of this project. Such settings would substantially increase
training and evaluation cost, particularly for repeated dynamic simulations.
Instead, the study focuses on Solomon-scale instances, which are widely used in
VRPTW benchmarking and are appropriate for the methodological objectives of this
work. In addition, the full canonical 100-customer Solomon instances were not
used for all AI training experiments, as training RouteFinder directly on the
complete instances proved computationally expensive within the available project
resources. To maintain feasible training times while still preserving meaningful
problem scale, experiments were conducted on reduced variants consisting of the
first 50 customers and the first 75 customers from each original instance. These
subsets allow controlled analysis of scaling behaviour while retaining the
underlying spatial and temporal structure of the benchmark instances.
    \item \textbf{Customer rejection is immediate.}
A newly revealed request was treated as rejected if the current re-optimization
step failed to incorporate it. This was adopted as a deliberate operational
assumption to model settings in which requests must be accepted or rejected at
the moment of disclosure, rather than kept pending for later insertion. The
choice preserves a clear and consistent decision rule across methods, but it
also defines a stricter service policy than formulations that allow deferred
acceptance. Consequently, rejection counts should be interpreted within the
context of this immediate-decision setting.

\end{itemize}

These limitations should be considered when interpreting the results. The
findings are most directly applicable to controlled dynamic VRPTW settings with
deterministic travel conditions and medium-scale benchmark instances. Extending
the framework to richer operational environments remains a natural direction for
future work.
